\hypertarget{appendix-7107009c505b}{}
\section*{Local rigidity and mesh orientation fields}
\chapter{Local rigidity and mesh orientation fields}\label{appendix-7107009c505b-local-rigidity-and-mesh-orientation-fields}

\textbf{Curriculum depth:} Developed · \textbf{First seen:} University courses · \textbf{Latest development:} University courses

\section{Definition and mechanism}\label{appendix-7107009c505b-definition-and-mechanism}

As-Rigid-As-Possible (ARAP) deformation seeks vertex positions that change as little as possible from a preferred local rigid transform. For each vertex neighborhood, fit a rotation to the deformed edge vectors, then minimize residuals between transformed rest edges and current edges. Alternating between local rotations and a global position solve makes a nonlinear shape-editing problem tractable, but the result depends on constraints, initialization, mesh quality, and rotation fitting. Instant Meshes uses direction or orientation fields over a surface to guide a quad-dominant remeshing; field singularities encode where a globally consistent direction cannot exist. A mesh's half-edge connectivity, boundary structure, and orientation determine valid neighborhoods and field transport. Poor triangles, seams, boundary conditions, or inconsistent orientation can create artifacts. ARAP is a local preservation objective, not exact rigidity of the whole shape; a direction field guides topology but does not uniquely determine a good mesh without sizing and feature constraints.

\section{Prerequisites}\label{appendix-7107009c505b-prerequisites}

\begin{itemize}
\tightlist
\item
  \hyperref[appendix-1aea26b2307f]{Half-edge mesh connectivity}
\item
  \href{../04-Visual-Computing-II/main.pdf\#nameddest=appendix-fab348609369}{Manifold polygon meshes and boundary structure (04-Visual-Computing-II, appendix-fab348609369)}
\item
  \href{../04-Visual-Computing-II/main.pdf\#nameddest=appendix-0234ab7f86b1}{Subdivision surface masks (04-Visual-Computing-II, appendix-0234ab7f86b1)}
\item
  \href{../04-Visual-Computing-II/main.pdf\#nameddest=appendix-29f4bea4c2bb}{Constrained optimization and inverse kinematics (04-Visual-Computing-II, appendix-29f4bea4c2bb)}
\item
  \href{../02-Foundations-II/main.pdf\#nameddest=appendix-2391db8052b1}{Cross products and oriented area (02-Foundations-II, appendix-2391db8052b1)}
\end{itemize}

\section{Used by}\label{appendix-7107009c505b-used-by}

\begin{itemize}
\tightlist
\item
  \href{Continuous\%20Notes/University\%20-\%20Worlds,\%20Images\%20and\%20Inference\%20-\%20Graphics,\%20Vision\%20and\%20Photography\#Meshes,\%20connectivity,\%20and\%20geometry\%20processing}{Meshes, connectivity, and geometry processing}
\end{itemize}

\section{Full teaching explanation}\label{appendix-7107009c505b-full-teaching-explanation}

\begin{itemize}
\tightlist
\item
  \href{Continuous\%20Notes/University\%20-\%20Worlds,\%20Images\%20and\%20Inference\%20-\%20Graphics,\%20Vision\%20and\%20Photography\#Meshes,\%20connectivity,\%20and\%20geometry\%20processing}{Meshes, connectivity, and geometry processing}
\end{itemize}

\section{Recall and next depth}\label{appendix-7107009c505b-recall-and-next-depth}

\begin{itemize}
\tightlist
\item
  Describe one local rotation/global solve step of ARAP and the role of a field singularity in remeshing.
\item
  Why can low local deformation energy fail to ensure a globally rigid shape?
\end{itemize}

\section{Source and navigation}\label{appendix-7107009c505b-source-and-navigation}

Teaching narrative: \href{Continuous\%20Notes/University\%20-\%20Worlds,\%20Images\%20and\%20Inference\%20-\%20Graphics,\%20Vision\%20and\%20Photography\#Meshes,\%20connectivity,\%20and\%20geometry\%20processing}{open the main explanation}. Browse the \href{../02-Foundations-II/main.pdf\#nameddest=appendix-a6d556a4ef50-university-graphics-vision-and-photography---independent-series}{University graphics and vision index (02-Foundations-II, appendix-a6d556a4ef50-university-graphics-vision-and-photography-\/-\/-independent-series)}. Additional chapter connections: \href{Continuous\%20Notes/University\%20-\%20Worlds,\%20Images\%20and\%20Inference\%20-\%20Graphics,\%20Vision\%20and\%20Photography\#Meshes,\%20connectivity,\%20and\%20geometry\%20processing}{Meshes, connectivity, and geometry processing}; \href{Continuous\%20Notes/University\%20-\%20Worlds,\%20Images\%20and\%20Inference\%20-\%20Graphics,\%20Vision\%20and\%20Photography\#Surfaces,\%20patches,\%20topology,\%20and\%20parameterization}{Surfaces, patches, topology, and parameterization}; \href{Continuous\%20Notes/Computing\%20Foundations\%20-\%20From\%20Bits\%20to\%20C++\%20and\%20Systems\%20APIs\#Put\%20the\%20algorithms\%20in\%20a\%20graphics\%20context}{Put the algorithms in a graphics context}.

Caption/source records: \href{Continuous\%20Notes/University\%20-\%20Worlds,\%20Images\%20and\%20Inference\%20-\%20Graphics,\%20Vision\%20and\%20Photography\#P3-16\%20-\%20Subdivision\%20Surfaces\%20Part\%201}{P3-16}, \href{Continuous\%20Notes/University\%20-\%20Worlds,\%20Images\%20and\%20Inference\%20-\%20Graphics,\%20Vision\%20and\%20Photography\#P3-17\%20-\%20Subdivision\%20Surfaces\%20Part\%202}{P3-17}, \href{Continuous\%20Notes/University\%20-\%20Worlds,\%20Images\%20and\%20Inference\%20-\%20Graphics,\%20Vision\%20and\%20Photography\#SC-15\%20-\%20Stanford\%20CS231N\%20Deep\%20Learning\%20for\%20Computer\%20Vision\%20Spring\%202025\%20Lecture\%2015:\%203D\%20Vision}{SC-15}.
